\documentclass[../../main.tex]{subfiles}
\graphicspath{{\subfix{../../image/}}}

\begin{document}

\section{解的几何性质}\label{section:ch09-5-解的几何性质}

\subsection{星形水平集}\label{subsection:ch09-5-1-星形水平集}

本小节解释一个偶尔有用的简单方法，用于研究各种 PDE 解水平集的几何性质。最简单情形出现在考察开集 $U$ 中调和函数时，$U$ 具有形式
\begin{align*}
U=W\setminus\overline{V},
\end{align*}
其中 $V\subset\subset W$，开集 $V$、$W$ 各自关于 $0$ 星形。写
\begin{align*}
\Gamma_0=\partial W,\qquad \Gamma_1=\partial V.
\end{align*}
考察问题
\begin{align}
\begin{cases}
\Delta u=0&\text{在 }U\text{ 中},\\
u=1&\text{在 }\Gamma_1\text{ 上},\\
u=0&\text{在 }\Gamma_0\text{ 上}.
\end{cases}\label{align:ch09-5-eq01}
\end{align}
物理上，一旦固定 $\Gamma_1$ 上势值为 $1$、$\Gamma_0$ 上势值为 $0$，$u$ 对应于区域 $U$ 内产生的静电势。由强极值原理，$0<u<1$ 在 $U$ 内。

\begin{theorem}[星形水平集]\label{theorem:ch09-5-1-星形水平集}
对每个 $0<\lambda<1$，水平集
\begin{align*}
\Gamma_\lambda\triangleq\{x\in U\mid u(x)=\lambda\}
\end{align*}
是光滑曲面，且是一个关于 $0$ 星形集合的边界。
\begin{proof}
1. 对每个 $\mu>0$，函数 $x\mapsto u(\mu x)$ 调和，从而
\begin{align*}
v(x)\triangleq\left.\frac{\mathrm{d}}{\mathrm{d}\mu}(u(\mu x))\right|_{\mu=1}=Du(x)\cdot x\qquad (x\in U)
\end{align*}
也调和。现因 $u=0$ 在 $\Gamma_0$ 上，在每个点 $x\in\partial W$ 处 $Du(x)$ 指向 $-\boldsymbol{\nu}(x)$ 方向。此外，因 $W$ 关于 $0$ 星形，$x\cdot\boldsymbol{\nu}(x)\geqslant 0$ 在 $\Gamma_0$ 上。故 $v=Du\cdot x\leqslant 0$ 在 $\Gamma_0$ 上。类似地 $v\leqslant 0$ 在 $\Gamma_1$ 上。由调和函数强极值原理，$v<0$ 在 $U$ 中。特别地，$Du\neq 0$ 在 $U$ 内。故隐函数定理（\S C.7）蕴涵 $\Gamma_\lambda$ 对 $0<\lambda<1$ 是光滑曲面。
2. 把 $u$ 延拓为在整个 $V$ 上等于 $1$，写
\begin{align*}
U_\lambda\triangleq\{x\in W\mid u>\lambda\}.
\end{align*}
则 $U_\lambda$ 是 $W$ 的开子集，$\partial U_\lambda=\Gamma_\lambda$。由强极值原理，$U_\lambda$ 连通。
现令 $x\in\Gamma_\lambda$，$\boldsymbol{\nu}(x)$ 表示 $\Gamma_\lambda$ 在 $x$ 处的外单位法向量。则 $Du(x)$ 指向 $-\boldsymbol{\nu}(x)$ 方向。因 $v(x)<0$，有 $x\cdot\boldsymbol{\nu}(x)>0$。此不等式对每个 $x\in\Gamma_\lambda$ 成立。
3. 由此 $\Gamma_\lambda$ 是一个关于 $0$ 星形集合的边界。为看清此点，回到 \S9.4.2 引理的证明，注意若 $\Gamma_\lambda$ 不关于 $0$ 星形，则可找到点 $x\in\Gamma_\lambda$，使 $y=\mu x\notin U_\lambda$，当 $\mu$ 充分接近 $1$、$\mu<1$。但随后可导出矛盾
\begin{align*}
\frac{\boldsymbol{\nu}(x)\cdot x}{|x|}=-\lim_{\mu\to 1^-}\frac{\boldsymbol{\nu}(x)\cdot(\mu x-x)}{|\mu x-x|}\leqslant 0.
\end{align*}
\end{proof}
\end{theorem}

\subsection{径向对称性}\label{subsection:ch09-5-2-径向对称}

本小节取 $U=B_0(0,1)$ 为 $\mathbb{R}^n$ 中的开单位球，考察一个半线性 Poisson PDE 的边值问题：
\begin{align}
\begin{cases}
-\Delta u=f(u)&\text{在 }U\text{ 中},\\
u=0&\text{在 }\partial U\text{ 上}.
\end{cases}\label{align:ch09-5-eq02}
\end{align}
我们关心正解：
\begin{align}
u>0\quad\text{在 }U\text{ 中},\label{align:ch09-5-eq03}
\end{align}
并设 $f:\mathbb{R}\to\mathbb{R}$ Lipschitz 连续，此外任意。目的是证明 $u$ 必为径向的，即 $u(x)$ 只依赖于 $r=|x|$。这是一个出人意料地强的结论，因为我们本质上未对非线性项作任何假设。

\subsubsection*{a. 极值原理}\label{subsubsection:ch09-5-2-a-极值原理}

证明将依赖于二阶椭圆 PDE 极值原理的一个推广。

\begin{lemma}[Hopf 引理的一个精化]\label{lemma:ch09-5-1-Hopf精化}
设 $V\subset\mathbb{R}^n$ 开，$v\in C^2(\overline{V})$，$c\in L^\infty(V)$。假设
\begin{align}
\begin{cases}
-\Delta v+cv\geqslant 0&\text{在 }V\text{ 中},\\
v\geqslant 0&\text{在 }V\text{ 中}.
\end{cases}\label{align:ch09-5-eq04}
\end{align}
又设 $v\not\equiv 0$。
(i) 若 $x_0\in\partial V$，$v(x_0)=0$，且 $V$ 在 $x_0$ 处满足内球条件，则
\begin{align}
\frac{\partial v}{\partial\nu}(x_0)<0.\label{align:ch09-5-eq05}
\end{align}
(ii) 此外，
\begin{align}
v>0\quad\text{在 }V\text{ 中}.\label{align:ch09-5-eq06}
\end{align}
注意这里关于零阶系数 $c$ 的符号不作任何假设。
\begin{proof}
令 $w\triangleq e^{-\lambda x_1}v$，$\lambda>0$ 下面选取。则 $v=e^{\lambda x_1}w$，故
\begin{align*}
cv\geqslant\Delta v=\Delta(e^{\lambda x_1}w)=\lambda^2v+2\lambda e^{\lambda x_1}w_{x_1}+e^{\lambda x_1}\Delta w.
\end{align*}
于是若 $\lambda=\|c\|_{L^\infty}^{1/2}$，
\begin{align*}
-\Delta w-2\lambda w_{x_1}\geqslant(\lambda^2-c)w\geqslant 0\quad\text{在 }V\text{ 中}.
\end{align*}
故 $w$ 是椭圆算子 $Kw\triangleq-\Delta w-2\lambda w_{x_1}$（无零阶项）的上解。强极值原理蕴涵 $w>0$ 在 $V$ 中。由 Hopf 引理（\S6.4.2），$\frac{\partial w}{\partial\nu}(x_0)<0$。但因 $v(x_0)=0$，
\begin{align*}
\frac{\partial w}{\partial\nu}(x_0)=Dw(x_0)\cdot\boldsymbol{\nu}(x_0)=e^{-\lambda x_1^0}\frac{\partial v}{\partial\nu}(x_0).
\end{align*}
故断言 (i) 成立；断言 (ii) 因 $w>0$ 在 $V$ 中而成立。
\end{proof}
\end{lemma}

\begin{lemma}[边界估计]\label{lemma:ch09-5-2-边界估计}
设 $u\in C^2(\overline{U})$ 满足\eqref{align:ch09-5-eq02}、\eqref{align:ch09-5-eq03}。则对每个点 $x_0\in\partial U\cap\{x_n>0\}$，或者
\begin{align}
u_{x_n}(x_0)<0,\label{align:ch09-5-eq07}
\end{align}
或者
\begin{align}
u_{x_n}(x_0)=0,\quad u_{x_nx_n}(x_0)>0.\label{align:ch09-5-eq08}
\end{align}
无论哪种情形，$u$ 作为 $x_n$ 的函数在 $x_0$ 附近严格递减。
\begin{proof}
1. 固定任一点 $x_0\in\partial U\cap\{x_n>0\}$，令 $\nu=\boldsymbol{\nu}(x_0)=(\nu_1,\cdots,\nu_n)$ 表示 $\partial U$ 在 $x_0$ 处的外单位法向量。注意 $\nu_n>0$。
2. 先断言：若
\begin{align}
f(0)\geqslant 0,\label{align:ch09-5-eq09}
\end{align}
则 $u_{x_n}(x_0)<0$。事实上
\begin{align*}
0=-\Delta u-f(u)=-\Delta u-f(u)+f(0)-f(0)\leqslant-\Delta u+cu,
\end{align*}
其中 $c(x)\triangleq-\int_0^1 f'(su(x))\,\mathrm{d}s$。由\cref{lemma:ch09-5-1-Hopf精化}，$\frac{\partial u}{\partial\nu}(x_0)<0$。因 $Du$ 在 $\partial U$ 上平行于 $\boldsymbol{\nu}$ 且 $\nu_n>0$，得 $u_{x_n}(x_0)<0$。
3. 现假设
\begin{align}
f(0)<0.\label{align:ch09-5-eq10}
\end{align}
若 $u_{x_n}(x_0)<0$，已毕。否则，因 $Du$ 平行于 $\boldsymbol{\nu}$，
\begin{align}
Du(x_0)=0.\label{align:ch09-5-eq11}
\end{align}
因\eqref{align:ch09-5-eq02} 在坐标轴旋转下不变，不妨设 $x_0=(0,\cdots,1)$、$\nu=(0,\cdots,1)$。
4. 断言
\begin{align}
u_{x_ix_j}(x_0)=-f(0)\nu_i\nu_j\label{align:ch09-5-eq12}
\end{align}
对每个 $i,j=1,\cdots,n$。因 $u=0$ 在 $\partial U$ 上，对所有 $x'\in\mathbb{R}^{n-1}$、$|x'|\leqslant 1$，$u(x',\gamma(x'))=0$，其中 $\gamma(x')=(1-|x'|^2)^{1/2}$。关于 $x_i$、$x_j$（$i,j=1,\cdots,n-1$）求导并用\eqref{align:ch09-5-eq11}，得
\begin{align}
u_{x_ix_j}(x_0)=0\qquad (i,j=1,\cdots,n-1).\label{align:ch09-5-eq13}
\end{align}
因 $u_{x_n}\leqslant 0$ 在 $\partial U\cap\{x_n>0\}$ 上且 $u_{x_n}(x_0)=0$，映射 $x'\mapsto u_{x_n}(x',\gamma(x'))$ 在 $x'=0$ 处取最大值。故
\begin{align}
u_{x_nx_i}(x_0)=0\qquad (i=1,\cdots,n-1).\label{align:ch09-5-eq14}
\end{align}
最后，\eqref{align:ch09-5-eq13}、\eqref{align:ch09-5-eq14} 与 PDE \eqref{align:ch09-5-eq02} 迫使 $u_{x_nx_n}(x_0)=-f(0)$。此等式即 $\nu=(0,\cdots,1)$ 时的\eqref{align:ch09-5-eq12}。回到原坐标轴，得\eqref{align:ch09-5-eq12}。
5. 在\eqref{align:ch09-5-eq12} 中令 $i=j=n$，用\eqref{align:ch09-5-eq10} 得
\begin{align*}
u_{x_nx_n}(x_0)=-f(0)\nu_n\nu_n>0.
\end{align*}
\end{proof}
\end{lemma}

\subsubsection*{b. 移动平面}\label{subsubsection:ch09-5-2-b-移动平面}

接着引入一个"移动平面" $P_\lambda$，将跨此平面反射偏微分方程。

\paragraph{记号。}
(i) 若 $0\leqslant\lambda\leqslant 1$，定义平面
\begin{align*}
P_\lambda\triangleq\{x\in\mathbb{R}^n\mid x_n=\lambda\}.
\end{align*}
(ii) 写 $x_\lambda\triangleq(x_1,\cdots,x_{n-1},2\lambda-x_n)$ 表示 $x$ 关于 $P_\lambda$ 的反射。
(iii) $E_\lambda\triangleq\{x\in U\mid\lambda<x_n<1\}$。

\begin{figure}[htbp]
\centering
\includegraphics[width=0.6\textwidth]{ch09-fig-reflection-plane.png}
\caption{跨平面的反射}\label{figure:ch09-reflection-plane}
\end{figure}

\begin{theorem}[径向对称性]\label{theorem:ch09-5-2-径向对称}
设 $u\in C^2(\overline{U})$ 求解\eqref{align:ch09-5-eq02}、\eqref{align:ch09-5-eq03}。则 $u$ 是径向的；即
\begin{align*}
u(x)=v(r)\qquad (r=|x|)
\end{align*}
对某严格递减函数 $v:[0,1]\to[0,\infty)$。
\begin{proof}
1. 对每个 $0\leqslant\lambda<1$，考察命题
\begin{align}
(15_\lambda)\qquad u(x)<u(x_\lambda)\text{ 对每个点 }x\in E_\lambda.
\end{align}
2. 由\cref{lemma:ch09-5-2-边界估计}，$(15_\lambda)$ 对每个 $\lambda<1$、$\lambda$ 充分接近 $1$ 成立。令
\begin{align}
\lambda_0\triangleq\inf\{0\leqslant\lambda<1\mid(15_\mu)\text{ 对每个 }\lambda\leqslant\mu<1\text{ 成立}\}.\label{align:ch09-5-eq16}
\end{align}
将证明
\begin{align}
\lambda_0=0.\label{align:ch09-5-eq17}
\end{align}
反设 $\lambda_0>0$。写 $w(x)\triangleq u(x_{\lambda_0})-u(x)$（$x\in E_{\lambda_0}$）。则
\begin{align*}
-\Delta w=f(u(x_{\lambda_0}))-f(u(x))=-cw\quad\text{在 }E_{\lambda_0}\text{ 中},
\end{align*}
其中 $c(x)\triangleq-\int_0^1 f'(su(x_{\lambda_0})+(1-s)u(x))\,\mathrm{d}s$。因 $w\geqslant 0$ 在 $E_{\lambda_0}$ 中，由\cref{lemma:ch09-5-1-Hopf精化}（应用于 $V=E_{\lambda_0}$）推得 $w>0$ 在 $E_{\lambda_0}$ 中，$w_{x_n}>0$ 在 $P_{\lambda_0}\cap U$ 上。故
\begin{align}
u(x)<u(x_{\lambda_0})\quad\text{在 }E_{\lambda_0}\text{ 中},\label{align:ch09-5-eq18}
\end{align}
且
\begin{align}
u_{x_n}<0\quad\text{在 }P_{\lambda_0}\cap U\text{ 上}.\label{align:ch09-5-eq19}
\end{align}
用\eqref{align:ch09-5-eq18}、\eqref{align:ch09-5-eq19} 与\cref{lemma:ch09-5-2-边界估计}，得
\begin{align}
u(x)<u(x_{\lambda_0-\varepsilon})\quad\text{在 }E_{\lambda_0-\varepsilon}\text{ 中，对所有 }0\leqslant\varepsilon\leqslant\varepsilon_0,\label{align:ch09-5-eq20}
\end{align}
若 $\varepsilon_0$ 充分小。若 $\lambda_0>0$，断言\eqref{align:ch09-5-eq20} 与 $\lambda_0$ 的选取\eqref{align:ch09-5-eq16} 矛盾。
3. 因 $\lambda_0=0$，看到 $u(x_1,\cdots,x_{n-1},-x_n)\geqslant u(x_1,\cdots,x_n)$ 对所有 $x\in U\cap\{x_n>0\}$。在 $U\cap\{x_n<0\}$ 中类似论证表明 $u(x_1,\cdots,x_{n-1},-x_n)\leqslant u(x_1,\cdots,x_n)$ 对所有 $x\in U\cap\{x_n>0\}$。故 $u$ 关于平面 $P_0$ 对称，且 $u_{x_n}=0$ 在 $P_0$ 上。
此论证在坐标轴任意旋转后同样适用，故定理得证。
\end{proof}
\end{theorem}

\end{document}
